\documentclass[11pt,a4paper]{article}
\usepackage[margin=27mm]{geometry}
\usepackage{amsmath,amssymb,amsthm,mathtools,booktabs,array,longtable}
\usepackage[T1]{fontenc}
\usepackage{lmodern,microtype}
\usepackage{xcolor}
\usepackage[unicode,colorlinks=true,linkcolor=blue!45!black,citecolor=blue!45!black,urlcolor=blue!45!black]{hyperref}
\usepackage{enumitem}
\hypersetup{pdftitle={Exact Certificates for Unimodality of Forest Independence Polynomials},pdfauthor={Tong Zhang}}
\setlist{nosep}
\allowdisplaybreaks
\emergencystretch=2em
\newtheorem{theorem}{Theorem}[section]
\newtheorem{lemma}[theorem]{Lemma}
\newtheorem{proposition}[theorem]{Proposition}
\newtheorem{corollary}[theorem]{Corollary}
\theoremstyle{definition}\newtheorem{definition}[theorem]{Definition}
\theoremstyle{remark}\newtheorem{remark}[theorem]{Remark}
\newcommand{\E}{\mathbb E}
\newcommand{\Var}{\operatorname{Var}}
\newcommand{\Prb}{\mathbb P}
\newcommand{\ind}{\mathbf1}
\newcommand{\pos}[1]{(#1)_+}
\title{Exact Certificates for Unimodality of Forest Independence Polynomials\\[3mm]\large A consolidated proof manuscript and reproducible supplement}
\author{Tong Zhang\\\small\texttt{ZhangZenoZhang@gmail.com}}
\date{26 September 2026\quad Version 1.1 (revised submission draft)}
\begin{document}
\maketitle
\begin{abstract}
We assemble two certificate-assisted arguments for the independence sequences of finite simple forests. The finite-order argument covers every forest on at most 60 vertices by structural counting inequalities, 670 rational separating inequalities, and 1907 nonnegative rational linear certificates. A separate hard-core argument, supplied with a frozen dependency chain, asserts unimodality for all forests on at least 59 vertices. Their scopes overlap, so that no order is omitted. We give the finite-order structural proof, the conditional-mixture reduction and new matching-count constraints for the large-order argument, and a precise map to the full inherited mathematical and computational supplements. Acceptance uses integer or directed rational arithmetic; numerical optimizers are not part of the verification. This is a mathematical proof submission for review, not a claim of completed Lean formalization or prior journal acceptance. A rational obstruction at order 100 concerns only the finite-order relaxation and is not a graph counterexample.
\end{abstract}
\noindent\textbf{Keywords.} Forest; independent set; unimodality; hard-core model; exact certificate; computer-assisted proof.

\section{The statement, scope, and proof architecture}\label{sec:scope}
For a finite simple graph $F$, write
\[
 I_F(x)=\sum_{r=0}^{\alpha(F)}c_r(F)x^r,
 \qquad c_r(F)=|\{S\subseteq V(F): S\text{ independent},\ |S|=r\}|.
\]
A finite sequence is unimodal when it weakly increases to some index and weakly decreases thereafter. We include the empty forest, with polynomial 1, as well as isolated vertices and disconnected forests. Throughout, $n$ denotes the number of vertices; it never denotes the matching number or a coefficient index.

The tree and forest questions originate in Alavi, Malde, Schwenk, and Erd\H{o}s \cite{alavi} and appear as Erd\H{o}s Problem 993 and Justin Sun Prize Problem JSP-000826 \cite{erdos,jsp}. The forest assertion cannot be inferred merely by multiplying arbitrary unimodal tree polynomials: unimodality is not generally preserved by convolution. Our two inputs each quantify over forests directly.

\begin{theorem}[Consolidated assertion submitted for review]\label{thm:all}
The independence sequence of every finite simple forest is unimodal.
\end{theorem}
The proof is organized through the following two propositions.
\begin{proposition}[Finite-order part]\label{prop:finite}
Every finite simple forest on at most 60 vertices has a unimodal independence sequence.
\end{proposition}
\begin{proposition}[Large-order part]\label{prop:large}
Every finite simple forest on at least 59 vertices has a unimodal independence sequence.
\end{proposition}
\begin{proof}[Proof of Theorem~\ref{thm:all} from the two propositions]
If $n\le60$, use Proposition~\ref{prop:finite}. If $n\ge61$, use Proposition~\ref{prop:large}. The two propositions also independently cover $n=59,60$. No induction across these orders and no closure under disjoint union is needed.
\end{proof}

Proposition~\ref{prop:finite} is developed below with its exact certificate interface. Proposition~\ref{prop:large} uses the large-order manuscript and frozen dependency chain distributed as Supplement A. We reproduce its final extension and the principal analytic interface in Sections~\ref{sec:window}--\ref{sec:repro}. The full analytic proofs and earlier thresholds are integral parts of the supplement, not assumptions supplied by a numerical success flag. The historical wording ``conditional on inherited structural lemmas'' in the input releases refers to those stated dependencies. To accept the consolidated theorem, a reviewer must accept their proofs as well as the finite arithmetic; two overlapping numerical ranges alone are not a proof.

Fang, Lu, Nevo, Yao, and Zheng \cite{fang} prove unimodality for every sufficiently large forest. Their Theorem 1.1 states the existence of an absolute threshold; it does not supply the value 59 used here. We do not substitute that existence statement for the explicit threshold chain, nor claim priority over their asymptotic result. Counterexamples to log-concavity of tree independence sequences \cite{kadrawi,ramos} do not contradict the unimodality assertion. We make no assertion of log-concavity of the full sequence.

\paragraph{Reading the supplement.}
Supplement A is the original \nolinkurl{forest_threshold_59_audited_handoff.zip}; its nested archives preserve the stages $80,99,130,170,350,500,900$. Supplement B is \nolinkurl{forest_n60_extension_and_n100_gap.zip}. A separate Chinese finite-order proof supplied with the files is byte-identical to its copy in Supplement B. The distribution includes those original bytes, a flattened map of mathematical source documents, the present English manuscript, and fresh audit evidence. Supplementary Chinese texts retain their original notation and historical scope statements.

\section{A finite certificate theorem for all forests of order at most sixty}
\label{f60:sec-main}

For a finite simple graph $F$, let $c_r(F)$ denote the number of independent
sets of size $r$, and let
\[
 I_F(z)=\sum_{r=0}^{\alpha(F)}c_r(F)z^r
\]
be its independence polynomial.  We use the non-strict convention for
unimodality: a sequence is unimodal if it is nondecreasing up to some index
and nonincreasing thereafter.  Thus a plateau at the maximum is allowed.
The result in this section concerns arbitrary forests, including disconnected
forests and forests with isolated vertices.

\begin{theorem}[Finite-order forest unimodality]
\label{f60:thm-main}
If $F$ is a forest with at most $60$ vertices, then the coefficient sequence
of $I_F(z)$ is unimodal.
\end{theorem}

The proof has two parts.  Structural inequalities place the independent-set
counts of every forest in an explicitly defined real polyhedron.  A finite
collection of rational certificates then establishes the required signs of
successive coefficient differences throughout the relevant parameter
domain.  The certificate stage does not enumerate forests.  It is,
however, an essential part of the proof: the result is not asserted to have
a certificate-free closed-form proof.  Neither Theorem~\ref{f60:thm-main}
nor the argument below asserts log-concavity of the entire coefficient
sequence, unimodality for all orders, or a corresponding theorem up to
order $100$.

Throughout this section, binomial coefficients are defined by
\[
 \binom{b}{r}=0\qquad\text{unless }b,r\text{ are integers with }0\le r\le b.
\]
For polynomials with real coefficients, $P\le_{\mathrm{coeff}}Q$ means
that every coefficient of $Q-P$ is nonnegative.  We write
$x_+=\max(x,0)$ and use $\mathbf{1}_{A}$ for the indicator of a condition
$A$.

\subsection{A decomposition relative to a maximum independent set}
\label{f60:subsec-decomposition}

Let $F$ be a forest on $n$ vertices.  Write $a=\alpha(F)$ and let $v$ be
its matching number.  Leaf deletion gives
\begin{equation}
 a+v=n,\qquad \delta:=a-v\ge0.
 \label{f60:eq-matching}
\end{equation}
For completeness, if $u$ is a leaf with neighbour $w$, a maximum
independent set may be chosen to contain $u$: replace $w$ by $u$ if
necessary.  Consequently
$\alpha(F)=1+\alpha(F-\{u,w\})$.  A maximum matching may likewise be
chosen to contain $uw$, by replacing an edge incident with $w$ by $uw$
when necessary.  Its size is therefore $1+\nu(F-\{u,w\})$.  An isolated
vertex increases the independence number by one and leaves the matching
number unchanged.  These observations prove the first identity in
\eqref{f60:eq-matching} by induction.  The inequality follows from
$2v\le n$.

Fix a maximum independent set $I$ and put $Y=V(F)\setminus I$, so that
$|I|=a$ and $|Y|=v$.  For any fixed maximum matching, an independent set
contains at most one endpoint of each matched edge, and there are
$\delta=n-2v$ unmatched vertices.  Since $|I|=v+\delta$, equality holds
in all of these restrictions.  Hence $I$ contains every unmatched vertex
and exactly one endpoint of every matched edge.  In particular, the
matching saturates $Y$ and matches it into $I$.

For an independent set $J\subseteq Y$, define
\[
 j=|J|,\qquad m(J)=|I\setminus N_I(J)|.
\]
Let $t_{j,m}$ count the independent subsets $J$ of $Y$ for which these
parameters equal $j,m$.  The set $J\cup(I\setminus N_I(J))$ is independent,
so $j+m\le a$.  By convention $t_{0,a}=1$, and all other $t_{0,m}$ vanish.
Every independent set of $F$ is uniquely obtained by first choosing its
intersection $J$ with $Y$ and then choosing an arbitrary subset of
$I\setminus N_I(J)$.  Thus
\begin{equation}
 I_F(z)=(1+z)^a+
 \sum_{j=1}^{v}\sum_{m=0}^{a-j}t_{j,m}z^j(1+z)^m.
 \label{f60:eq-decomposition}
\end{equation}
In particular,
\begin{equation}
 c_r=\binom ar+
 \sum_{j=1}^{v}\sum_{m=0}^{a-j}\binom m{r-j}t_{j,m}.
 \label{f60:eq-coefficients}
\end{equation}
No maximality requirement is placed on $J$; all independent sets are
included in this identity.

Define
\begin{equation}
 T_j=\sum_m t_{j,m},\qquad M_j=\sum_m mt_{j,m},\qquad
 W_j=\binom vj.
 \label{f60:eq-layer-notation}
\end{equation}
The basic constraints are
\begin{equation}
 t_{j,m}\ge0,\qquad T_1=v,\qquad T_j\le W_j\quad(2\le j\le v).
 \label{f60:eq-basic-counts}
\end{equation}
All certificate arguments below remain valid when the $t_{j,m}$ are
arbitrary nonnegative real numbers satisfying the displayed constraints.
Integrality is used in deriving some structural inequalities for actual
forests, but is not imposed in the subsequent linear relaxation.

\subsection{Coefficient bounds and the first cross-layer inequalities}
\label{f60:subsec-old-rows}

\begin{lemma}[Path lower bound and matching upper bound]
\label{f60:lem-path-matching}
For every forest of order $n$ and independence number $a$, with $v=n-a$,
\begin{equation}
 \binom{n-r+1}{r}\le c_r\qquad(r\ge0),
 \label{f60:eq-path-bound}
\end{equation}
and
\begin{equation}
 I_F(z)\le_{\mathrm{coeff}}(1+z)^{a-v}(1+2z)^v.
 \label{f60:eq-matching-bound}
\end{equation}
\end{lemma}

\begin{proof}
The number of independent $r$-sets of the path on $n$ vertices is
$p_{n,r}=\binom{n-r+1}{r}$.  A leaf $u$ with neighbour $w$ gives
\[
 I_F=I_{F-u}+zI_{F-\{u,w\}},
\]
and $p_{n,r}=p_{n-1,r}+p_{n-2,r-1}$.  Induction therefore proves
\eqref{f60:eq-path-bound} when a leaf is present.  If $u$ is isolated,
$I_F=(1+z)I_{F-u}$; the induction bound and
$p_{n-1,r-1}\ge p_{n-2,r-1}$ give the same conclusion.  The empty and
one-vertex forests start the induction.  Coefficients beyond the degree
are handled by the binomial convention.

For the upper bound, $m\le a-j$ in
\eqref{f60:eq-decomposition}, and the layer mass is at most $\binom vj$.
Consequently
\begin{align*}
 I_F(z)&\le_{\mathrm{coeff}}
 \sum_{j=0}^{v}\binom vj z^j(1+z)^{a-j}\\
 &=(1+z)^{a-v}(1+2z)^v.
\end{align*}
\end{proof}

\begin{lemma}[Private-neighbour inequality]
\label{f60:lem-private}
For $2\le r\le v$ and every integer $h\ge0$,
\begin{equation}
 \sum_m\bigl((r-1)m+a-r+1-rh\bigr)_+t_{r,m}
 \le(v-r+1)\sum_m(m-h)_+t_{r-1,m}.
 \label{f60:eq-private}
\end{equation}
\end{lemma}

\begin{proof}
Fix an independent $r$-set $J\subseteq Y$, put $m=m(J)$ and $d=a-m$,
and consider the bipartite forest induced by $J\cup N_I(J)$.
Each vertex of $J$ has a neighbour in $I$, since otherwise $I$ would not
be maximum.  If this forest has $b\ge1$ components and $p$ vertices of
degree one on its $I$ side, its number of edges is $r+d-b$, while summing
degrees on that side gives at least $p+2(d-p)$.  Therefore
\[
 p\ge d-r+b\ge d-r+1.
\]
For $u\in J$, let $p_u$ be the number of its private neighbours in $I$,
that is, neighbours of $u$ adjacent to no other member of $J$.  Then
\begin{equation}
 m(J\setminus\{u\})=m+p_u,\qquad
 \sum_{u\in J}p_u\ge d-r+1.
 \label{f60:eq-private-release}
\end{equation}
It follows that
\[
 \sum_{u\in J}(m+p_u-h)_+
 \ge \left(r(m-h)+\sum_{u\in J}p_u\right)_+
 \ge\bigl((r-1)m+a-r+1-rh\bigr)_+.
\]
Sum over $J$.  Each independent $(r-1)$-set of $Y$ is obtained by
deleting one vertex from at most $v-r+1$ independent $r$-sets.  The
resulting reverse count proves \eqref{f60:eq-private}.
\end{proof}

\begin{lemma}[Matching-based extension inequality]
\label{f60:lem-hall}
Put
\[
 H_{r,l}=\sum_m\binom ml t_{r,m},\qquad H_{0,l}=\binom al.
\]
Whenever the factor on the right is nonnegative,
\begin{equation}
 (r+1)H_{r+1,l}
 \le\bigl(v-r-\max(0,l-\delta)\bigr)H_{r,l}.
 \label{f60:eq-hall}
\end{equation}
\end{lemma}

\begin{proof}
Fix a set $L\subseteq I$ of size $l$.  At most $\delta$ vertices of $I$
are unmatched by a fixed maximum matching.  The matched vertices of $L$
therefore exclude at least $\max(0,l-\delta)$ distinct vertices of $Y$.
Among the remaining vertices, an independent $r$-set has at most
$v-r-\max(0,l-\delta)$ extensions.  Double-count the pairs consisting of
such an $r$-set and one extending vertex, and then sum over $L$.
An independent set $J$ contributes $\binom{m(J)}l$ choices of $L$,
giving precisely \eqref{f60:eq-hall}.
\end{proof}

\subsection{Choosing a maximum independent set with a sparse complement}
\label{f60:subsec-sparse}

The preceding formulas hold for every maximum independent set.  The
remaining inequalities are used with a single maximum independent set
chosen as follows.

\begin{lemma}[Sparse-complement selection]
\label{f60:lem-sparse}
A forest has a maximum independent set $I$ whose complement $Y$ satisfies
\begin{equation}
 e(F[Y])\le E:=\max(0,\delta-1).
 \label{f60:eq-sparse}
\end{equation}
If $v=|Y|\ge1$, then also $e(F[Y])\le v-1$.
\end{lemma}

\begin{proof}
Fix a maximum matching and a bipartite colouring of each component of
$F$.  Contract each matched edge to one paired node, and retain each
unmatched vertex as a terminal.  The contracted graph is a forest.
Every maximum independent set contains all terminals and selects exactly
one endpoint from each matched pair.  Label a paired node by the original
bipartite colour of its selected endpoint, and label a terminal by its own
colour.

Consider an edge of the contracted forest.  Its corresponding original
edge joins opposite bipartite colours.  If the two node labels agree, the
original edge has exactly one endpoint in $I$.  If they disagree, it has
either both endpoints in $I$ or both in $Y$; the former is excluded by
independence.  Thus edges between different labels correspond precisely
to edges of $F[Y]$.  No such edge is incident with a terminal.

Choose a maximum independent set minimizing $e(F[Y])$.  Suppose a
monochromatic connected block in the contracted forest has no terminal
and has a boundary edge.  Swap the selected and unselected endpoints in
every matched pair belonging to this block.  Internal edges still have
exactly one selected endpoint.  Every boundary edge, previously an edge
with both endpoints in $Y$, now has exactly one selected endpoint.
Matched edges still contribute one selected endpoint each.  Hence the
new set is independent, has the same size, and has strictly fewer edges
in its complement, a contradiction.

In an original component containing $\delta_c>0$ terminals, every
monochromatic block consequently contains a terminal.  There are at most
$\delta_c$ such blocks, and contracting them gives a tree, so there are
at most $\delta_c-1$ edges between blocks.  In a component without
terminals there can be only one monochromatic block, hence no such
edges.  Summing over components proves
\eqref{f60:eq-sparse}.  The final assertion follows because $F[Y]$ is
itself a forest.
\end{proof}

From now on, fix the set supplied by Lemma~\ref{f60:lem-sparse}.  Define
\begin{equation}
 D_0=\min(E,v-1),\qquad e=e(F[Y])=\binom v2-T_2.
 \label{f60:eq-actual-edges}
\end{equation}
The last identity is exact: a two-element subset of $Y$ fails to be
independent precisely when it is an edge.  All constraints below refer
to this same $I$, this same $Y$, and this same value of $e$.

\subsection{Two-sided information from deletion of a vertex}
\label{f60:subsec-two-sided}

For an independent $r$-set $J\subseteq Y$, retain the notation
$m=m(J)$, $d=a-m$, and put $L=d-r+1$.  Since $r+m\le a$, we have
$L\ge1$.  Moreover, applying maximality of the independence number to
$J\setminus\{u\}$ gives
$(r-1)+(m+p_u)\le a$.  Together with
\eqref{f60:eq-private-release} and the fact that private neighbours are
distinct, this yields
\begin{equation}
 0\le p_u\le L,\qquad L\le\sum_{u\in J}p_u\le d.
 \label{f60:eq-release-range}
\end{equation}

For $2\le r\le v$ and an integer $h\ge0$, define
\begin{equation}
 b_{r,h}(m)=
 \begin{cases}
 r,&m\ge h,\\
 1,&m<h\text{ and }rh\le a+(r-1)m,\\
 0,&\text{otherwise}.
 \end{cases}
 \label{f60:eq-step-function}
\end{equation}

\begin{lemma}[Lower deletion threshold]
\label{f60:lem-tail}
For every $2\le r\le v$ and integer $h\ge0$,
\begin{equation}
 \sum_m b_{r,h}(m)t_{r,m}
 \le(v-r+1)\sum_{m\ge h}t_{r-1,m}.
 \label{f60:eq-tail}
\end{equation}
\end{lemma}

\begin{proof}
If $m\ge h$, all $r$ deletions from $J$ give an available-vertex count
at least $h$.  Otherwise let $t=h-m\ge1$.  The condition in the second
line of \eqref{f60:eq-step-function} is $rt\le d$.  If every deletion
had $p_u<t$, integrality would imply
\[
 \sum p_u\le r(t-1)\le d-r<L,
\]
contrary to \eqref{f60:eq-release-range}.  Thus at least one deletion
has the required count.  Summing this lower bound over $J$ and using at
most $v-r+1$ extensions for each parent set proves the inequality.
\end{proof}

For an independent $(r-1)$-set $A\subseteq Y$, the number of unselected
vertices forbidden by its edges is at most the total number $e$ of edges
of $F[Y]$.  Hence $A$ has at least
\begin{equation}
 s_r=\max(0,v-r+1-E)
 \label{f60:eq-extension-lower}
\end{equation}
extensions to an independent $r$-set.

To bound deletion sums in the other direction, write
\[
 d=qL+\rho,\qquad 0\le\rho<L.
\]
The identity
$rL-d=(r-1)(d-r)\ge0$ gives $q\le r$; if $q=r$, then $\rho=0$.
Define
\begin{equation}
 U_{r,h}(m)=q(m+L-h)_+
 +\begin{cases}
 (m+\rho-h)_++(r-q-1)(m-h)_+,&q<r,\\
 0,&q=r,
 \end{cases}
 \label{f60:eq-U}
\end{equation}
and
\begin{equation}
 V_{r,h}(m)=
 \begin{cases}
 r,&h\le m,\\
 \min\bigl(r,\lfloor d/(h-m)\rfloor\bigr),&m<h\le m+L,\\
 0,&h>m+L.
 \end{cases}
 \label{f60:eq-V}
\end{equation}

\begin{lemma}[Upper deletion bounds]
\label{f60:lem-release-upper}
For $2\le r\le v$ and integer $h\ge0$,
\begin{align}
 s_r\sum_m(m-h)_+t_{r-1,m}
 &\le\sum_m U_{r,h}(m)t_{r,m},
 \label{f60:eq-release-upper}\\
 s_r\sum_{m\ge h}t_{r-1,m}
 &\le\sum_m V_{r,h}(m)t_{r,m}.
 \label{f60:eq-tail-upper}
\end{align}
\end{lemma}

\begin{proof}
For fixed $J$, consider the feasible release vector $(p_u)_{u\in J}$
in \eqref{f60:eq-release-range}.  The function
$p\mapsto(m+p-h)_+$ is convex and nondecreasing.  Under the upper
constraints $0\le p_u\le L$ and $\sum p_u\le d$, its sum is maximized
by concentrating mass: $q$ coordinates equal $L$, at most one further
coordinate equals $\rho$, and the others are zero.  This follows, for
example, by repeatedly transferring mass between two coordinates
strictly inside $[0,L]$ toward an endpoint, choosing a direction that
does not decrease the convex sum.  Monotonicity permits total mass $d$.
The resulting maximum is \eqref{f60:eq-U}.

For the indicator sum, all deletions qualify when $h\le m$.  If
$0<h-m\le L$, each qualifying coordinate consumes at least $h-m$ of
the total mass at most $d$, giving the middle bound in
\eqref{f60:eq-V}.  If $h-m>L$, none qualifies.

Now sum either deletion statistic over all independent $r$-sets.  The
upper bounds just obtained give the right sides of
\eqref{f60:eq-release-upper} and \eqref{f60:eq-tail-upper}.  Counting
the same incidence relation from its $(r-1)$-set side gives the left
sides, because every such set has at least $s_r$ extensions.  When
$s_r=0$, both inequalities remain valid and are automatic consequences
of nonnegativity.
\end{proof}

\subsection{Available-vertex means and the complement's actual edge count}
\label{f60:subsec-mean}

\begin{lemma}[Coverage union bound]
\label{f60:lem-union}
For $2\le r\le v$,
\begin{equation}
 \sum_m(a-m)t_{r,m}
 \le\binom{v-1}{r-1}\sum_m(a-m)t_{1,m}.
 \label{f60:eq-union}
\end{equation}
\end{lemma}

\begin{proof}
For each independent $r$-set $J\subseteq Y$, its covered vertices in
$I$ form a union of the neighbourhoods of its members.  Its cardinality
is at most the sum of those neighbourhood sizes.  A fixed vertex of $Y$
belongs to at most $\binom{v-1}{r-1}$ such sets.  Summing the union bound
therefore gives the stated inequality.
\end{proof}

For $2\le r\le v$, put
\begin{equation}
 \beta_r=\binom{v-2}{r-1}-(a-2)\binom{v-2}{r-2},\qquad
 C_r^0=a\binom{v-2}{r}+(\delta+1)\binom{v-2}{r-1}.
 \label{f60:eq-mean-constants}
\end{equation}

\begin{lemma}[Mean bound retaining the actual edge count]
\label{f60:lem-mean}
With $e=e(F[Y])$ as in \eqref{f60:eq-actual-edges},
\begin{equation}
 M_r\ge C_r^0+\beta_r e\qquad(2\le r\le v).
 \label{f60:eq-mean}
\end{equation}
\end{lemma}

\begin{proof}
For $i\in I$, let $d_i=|N(i)\cap Y|$.  For integer $0\le d\le v$,
the function $d\mapsto\binom{v-d}{r}$ is discretely convex: its forward
differences are $-\binom{v-d-1}{r-1}$ and are nondecreasing.  The line
through its values at $d=1,2$ is therefore a lower support line on the
integer domain.  Explicitly,
\[
 \binom{v-d_i}{r}\ge
 \binom{v-2}{r}+(2-d_i)\binom{v-2}{r-1}.
\]
The number of edges between $I$ and $Y$ is at most $n-1-e$.
Consequently
\begin{align*}
 \sum_{i\in I}\binom{v-d_i}{r}
 &\ge a\binom{v-2}{r}
   +(2a-n+1+e)\binom{v-2}{r-1}\\
 &=C_r^0+e\binom{v-2}{r-1}.
\end{align*}
The sum on the left counts pairs consisting of any $r$-subset of $Y$
and an available vertex of $I$, without yet requiring the subset to be
independent.

For an edge $uw$ of $F[Y]$, both $u$ and $w$ have nonempty
neighbourhoods in $I$, and those neighbourhoods are disjoint, since a
common neighbour would produce a triangle.  Every $r$-set containing
$u,w$ therefore has at most $a-2$ available vertices of $I$.  There are
$\binom{v-2}{r-2}$ candidate sets containing this edge.  A union bound
over the $e$ edges shows that deleting all nonindependent subsets from
the preceding count loses at most
$e(a-2)\binom{v-2}{r-2}$.  Subtracting this amount proves the lemma.
\end{proof}

The coefficient $\beta_r$ may be negative.  Accordingly we do not
replace $e$ independently by an endpoint of its allowed interval for
each $r$.  We retain the common variable and then substitute the exact
identity $e=\binom v2-T_2$; this is the source of the linear row used in
the certificates.

\subsection{Coefficient bounds for a forest with a fixed edge count}
\label{f60:subsec-edge-bounds}

\begin{lemma}[Edge-count coefficient bounds]
\label{f60:lem-edge-bounds}
For a forest $G$ on $v\ge1$ vertices and with $e$ edges,
\begin{equation}
 i_r(G)\ge\binom vr-e\binom{v-2}{r-2},
 \label{f60:eq-edge-lower-raw}
\end{equation}
where $i_r(G)$ counts its independent $r$-sets, and
\begin{equation}
 I_G(z)\le_{\mathrm{coeff}}
 (1+z)^{v-1}+z(1+z)^{v-e-1}.
 \label{f60:eq-star-upper}
\end{equation}
\end{lemma}

\begin{proof}
The first bound deletes from the $\binom vr$ subsets all those
containing an edge.  Each edge occurs in $\binom{v-2}{r-2}$ subsets,
so the union bound proves the claim, including cases when its right
side is negative.

For the upper bound, write
$B(v,e)=(1+z)^{v-1}+z(1+z)^{v-e-1}$.
If $e=0$, equality holds; orders at most two are immediate.  For
$v\ge3$ and $e>0$, choose a leaf $u$ with neighbour $w$ of degree $d$.
The two smaller forests have parameters $(v-1,e-1)$ and $(v-2,e-d)$.
Induction and leaf deletion give
\[
 I_G(z)\le_{\mathrm{coeff}}B(v-1,e-1)+zB(v-2,e-d).
\]
A direct subtraction gives
\[
 B(v,e)-B(v-1,e-1)-zB(v-2,e-d)
 =z^2\bigl((1+z)^{v-3}-(1+z)^{v-e+d-3}\bigr).
\]
The remaining forest on $v-2\ge1$ vertices has $e-d\le v-3$ edges,
so the second exponent is nonnegative.  Since $e\ge d$, it does not
exceed the first exponent.  The displayed difference has nonnegative
coefficients, completing the induction.
\end{proof}

Apply the lemma to $G=F[Y]$, whose independent-set coefficients are
$T_r$.  From \eqref{f60:eq-star-upper},
\begin{equation}
 T_r\le\binom{v-1}{r}+\binom{v-e-1}{r-1}.
 \label{f60:eq-edge-upper-raw}
\end{equation}
For $2\le r\le v$, the function on integer $0\le e\le v-1$ given by
\[
 g_r(e)=\binom{v-1}{r-1}-\binom{v-e-1}{r-1}
\]
is concave, satisfies $g_r(0)=0$, and has successive increments
$\binom{v-e-2}{r-2}$, which are nonincreasing.  If $D_0\ge1$, define
\begin{equation}
 G_r=\binom{v-1}{r-1}-\binom{v-D_0-1}{r-1}.
 \label{f60:eq-G}
\end{equation}
The chord bound gives $g_r(e)\ge eG_r/D_0$ for $0\le e\le D_0$.
Substituting \eqref{f60:eq-actual-edges} in the lower bound and in this
linearized upper bound yields
\begin{align}
 -T_r+\binom{v-2}{r-2}T_2
 &\le-\binom vr+\binom{v-2}{r-2}\binom v2,
 \label{f60:eq-edge-lower}\\
 D_0T_r-G_rT_2
 &\le D_0\binom vr-G_r\binom v2.
 \label{f60:eq-edge-upper}
\end{align}
When $D_0=0$, the implemented replacement for
\eqref{f60:eq-edge-upper} is the ordinary bound $T_r\le\binom vr$.
This replacement is weaker than all information available in that
case, but remains valid and is sufficient for the finite certificates.

\subsection{Reduction to successive differences on a finite parameter domain}
\label{f60:subsec-differences}

Set
\begin{equation}
 \Delta_k=c_k-c_{k-1},\qquad
 B_{m,j}(k)=\binom m{k-j}-\binom m{k-j-1}.
 \label{f60:eq-difference-notation}
\end{equation}
The exact coefficient identity gives
\begin{equation}
 \Delta_k=B_{a,0}(k)+\sum_{j,m}B_{m,j}(k)t_{j,m}.
 \label{f60:eq-difference}
\end{equation}

\begin{lemma}[Universal decreasing tail]
\label{f60:lem-tail-cutoff}
For a forest of independence number $a$,
\begin{equation}
 (k+1)c_{k+1}\le2(a-k)c_k.
 \label{f60:eq-tail-ratio}
\end{equation}
In particular, with
\begin{equation}
 K(a)=\left\lfloor\frac{2a+1}{3}\right\rfloor,
 \label{f60:eq-cutoff}
\end{equation}
the sequence $c_{K(a)},c_{K(a)+1},\ldots,c_a$ is nonincreasing.
\end{lemma}

\begin{proof}
Given an independent $k$-set $S$, let $A(S)$ be the vertices that can
be added to it.  The induced forest on $A(S)$ has independence number
at most $a-k$, since any of its independent sets can be joined to $S$.
A bipartite graph has an independent set containing at least half its
vertices, so $|A(S)|\le2(a-k)$.  Count pairs $(S,u)$ with $u\in A(S)$.
Every independent $(k+1)$-set occurs $k+1$ times, proving
\eqref{f60:eq-tail-ratio}.  The factor $2(a-k)/(k+1)$ is at most one
when $k\ge\lceil(2a-1)/3\rceil=K(a)$.
\end{proof}

For each $1\le k<K(a)$ it is sufficient to prove one of the statements
\begin{equation}
 \Delta_k>0,\qquad
 \Delta_{k+1}\le0,\qquad
 \Delta_k\le0\ \Longrightarrow\ \Delta_{k+1}\le0.
 \label{f60:eq-alternative}
\end{equation}
Indeed, once a nonpositive difference occurs, these alternatives
prevent a later positive difference before the tail.  The first
alternative cannot apply at an index where the difference is already
nonpositive, while the other two give the required next sign.
Lemma~\ref{f60:lem-tail-cutoff} supplies all remaining signs.  Also
$c_0=1$ and $c_1=n$.

The empty forest and forests with no edges have independence polynomial
$1$ or $(1+z)^n$, respectively, and are handled directly.  For forests
with at least one edge and order at most $60$, all required triples lie
in the explicitly regenerated domain
\begin{equation}
 \mathcal P_{60}=\left\{(n,a,k):\quad
 \begin{aligned}
 &2\le n\le60,\quad \lceil n/2\rceil\le a\le n-1,\\[-2pt]
 &1\le k<\lfloor(2a+1)/3\rfloor
 \end{aligned}\right\}.
 \label{f60:eq-domain}
\end{equation}
It contains $17\,100$ triples.  Empty ranges for small values of $a$
require no certificate.  These triples are scalar parameter choices,
not representatives of graph isomorphism classes.

\subsection{Exact certificate identities and the complete finite check}
\label{f60:subsec-certificates}

\paragraph{Direct criteria.}
Let
\begin{equation}
 q_{a,v}(r)=\sum_{j=0}^{\min(v,r)}\binom vj2^j\binom{a-v}{r-j},
 \label{f60:eq-q}
\end{equation}
the coefficient of $z^r$ in the upper-bound polynomial of
Lemma~\ref{f60:lem-path-matching}.  The comparisons
\[
 \binom{n-k+1}{k}>q_{a,v}(k-1)
 \quad\text{or}\quad
 q_{a,v}(k+1)\le\binom{n-k+1}{k}
\]
prove, respectively, $\Delta_k>0$ or $\Delta_{k+1}\le0$.

The basic layer masses give further direct bounds.  For each $j$ define
\[
 \begin{split}
 L_j(r)&=\begin{cases}
 \min_{0\le m\le a-1}B_{m,1}(r),&j=1,\\
 \min\bigl(0,\min_{0\le m\le a-j}B_{m,j}(r)\bigr),&j\ge2,
 \end{cases}\\
 U_j(r)&=\begin{cases}
 \max_{0\le m\le a-1}B_{m,1}(r),&j=1,\\
 \max\bigl(0,\max_{0\le m\le a-j}B_{m,j}(r)\bigr),&j\ge2.
 \end{cases}
 \end{split}
\]
Then
\begin{equation}
 B_{a,0}(r)+\sum_{j=1}^{v}W_jL_j(r)
 \le\Delta_r\le
 B_{a,0}(r)+\sum_{j=1}^{v}W_jU_j(r).
 \label{f60:eq-direct-layer}
\end{equation}
The absence of the extra zero in the first layer is justified by
$T_1=v$ exactly.  A strictly positive lower bound at $r=k$ or a
nonpositive upper bound at $r=k+1$ proves the corresponding alternative
in \eqref{f60:eq-alternative}.

For a rational number $s\ge0$, the same upper-bound rule can be applied
to $B_{m,j}(k+1)-sB_{m,j}(k)$.  A nonpositive resulting bound proves
\begin{equation}
 \Delta_{k+1}\le s\Delta_k,
 \label{f60:eq-separator}
\end{equation}
which proves the implication in \eqref{f60:eq-alternative}.  These are
the one-parameter rational comparison certificates called
\texttt{separator} in the data.

\paragraph{General linear certificates.}
Normalize the variables by $x_{j,m}=t_{j,m}/W_j$.  Their layer masses
are at most one, and the first layer mass is exactly one.
For every valid integer inequality
\[
 \sum_{j,m}a^{\rho}_{j,m}t_{j,m}\le b^{\rho},
\]
define its scale by
\begin{equation}
 N_\rho=\max\left(1,|b^{\rho}|,
       \max_{j,m}|W_ja^{\rho}_{j,m}|\right),
 \label{f60:eq-row-scale}
\end{equation}
and write the corresponding normalized row as
$A^\rho x\le b_\rho$, where
$A^\rho_{j,m}=W_ja^\rho_{j,m}/N_\rho$ and
$b_\rho=b^\rho/N_\rho$.
Normalize the target affine function in the same way, obtaining
\[
 Q(x)=q_0+\sum_{j,m}q_{j,m}x_{j,m}.
\]
If $\lambda_\rho\ge0$ and $\eta\in\mathbb R$, set
\[
 R_{j,m}=q_{j,m}-\sum_\rho\lambda_\rho A^\rho_{j,m}
                    -\eta\mathbf{1}_{j=1}.
\]
The following elementary identity is the entire certificate principle:
\begin{equation}
 Q(x)\le q_0+\sum_\rho\lambda_\rho b_\rho+\eta
       +\sum_{j=1}^{v}\max\bigl(0,\max_m R_{j,m}\bigr).
 \label{f60:eq-certificate}
\end{equation}
To see this, subtract the nonnegative combination of valid rows and use
the first-layer equality for the $\eta$ term.  On each remaining layer,
nonnegativity and mass at most one bound its residual contribution by
the nonnegative maximum of its residual coefficients.  No assertion
about linear-programming optimality is needed.

There are three linear certificate types.  A \texttt{positive}
certificate uses the target $-\Delta_k$ and makes the right side of
\eqref{f60:eq-certificate} strictly negative.  A \texttt{negative}
certificate uses $\Delta_{k+1}$ and makes the bound nonpositive.  A
\texttt{conditional} certificate has the latter target and is permitted
to add the row $\Delta_k\le0$.  It therefore proves exactly the
implication in \eqref{f60:eq-alternative}.  That assumption cannot be
discarded while retaining an unconditional negative target claim.

\paragraph{Exact arithmetic and exhaustive parameter coverage.}
Each stored linear certificate gives a positive integer denominator
$D_c$, integer row-weight numerators, and an integer numerator for the
unrestricted equality weight.  The principal checker
\texttt{verify.py} independently constructs the unweighted integer
rows, computes a common multiple of their normalization denominators,
and checks \eqref{f60:eq-certificate} using integers.  It regenerates
every triple in \eqref{f60:eq-domain}; it rejects missing, duplicate, or
extraneous triples, invalid row parameters, negative inequality
multipliers, and conditional assumptions with the wrong index.
Direct and one-parameter comparisons are likewise performed exactly.

A second arithmetic implementation, \texttt{verify\_grid.py}, uses
the weighted rows from \texttt{kernel.py}.  It rounds target coefficients
and right sides upward and row coefficients downward on an integer
grid.  Nonnegative variables and nonnegative row multipliers make this
an outward, conservative bound.  This second implementation checks all
$1\,907$ linear certificates; complete parameter coverage and the
$670$ one-parameter comparisons are checked by the principal verifier.
Both implementations use the Python standard library, without a graph
catalogue, an optimizer, or floating-point feasibility tolerances.

The complete partition obtained by rerunning these checks is shown in
Table~\ref{f60:tab-coverage}.  The matching-upper/path-lower test for a
nonpositive next difference is available in the verifier, but accounts
for no additional cases in this partition.

\begin{table}[htbp]
\centering
\begin{tabular}{lr}
\hline
Certificate or direct criterion & Number of triples\\
\hline
Path lower bound exceeds the preceding matching upper bound & $6\,724$\\
Direct positive lower bound for $\Delta_k$ & $5\,709$\\
Direct nonpositive upper bound for $\Delta_{k+1}$ & $2\,090$\\
One-parameter rational comparison & $670$\\
\texttt{positive} linear certificate & $603$\\
\texttt{negative} linear certificate & $158$\\
\texttt{conditional} linear certificate & $1\,146$\\
\hline
Total & $17\,100$\\
\hline
\end{tabular}
\caption{Complete parameter coverage for Theorem~\ref{f60:thm-main}.}
\label{f60:tab-coverage}
\end{table}

The data contain $2\,577$ stored triples: $670$ comparison parameters
and $1\,907$ linear certificates.  The latter contain $142\,017$
nonzero row-weight terms.  There are at most $610$ variables $t_{j,m}$
for any parameter pair in this domain; their values are not enumerated.
The largest common denominator used by the principal arithmetic check
has $68$ decimal digits.  The minimum verified normalized linear
margin is
\[
 \frac{953977}{42042000000}>0,
\]
attained at $(n,a,k)=(34,19,10)$.  This is a certificate upper-bound
margin after normalization, not a lower bound for an unnormalized
coefficient curvature.

The exact data file \texttt{certificates.json} is identified by the
SHA-256 digest
\begin{center}\small\ttfamily
a77f4fc43aec1b8eecab91c4418b6990\allowbreak
f72dee69d802129c78efd67d2e1bcc5a.
\end{center}
The archived proof data and checker source are mathematical supplements
to the finite theorem; the execution logs alone are not substitutes for
them.

\begin{proof}[Proof of Theorem~\ref{f60:thm-main}]
Forests without edges were dealt with above.  For a forest with an edge,
choose $I$ by Lemma~\ref{f60:lem-sparse} and form its exact counts
$t_{j,m}$.  The preceding lemmas show that these counts satisfy every
structural inequality used by the certificates.  Its parameters lie in
the stated ranges, and the exact finite check proves
\eqref{f60:eq-alternative} for every required index.  The first
nonpositive difference, if one occurs before $K(a)$, is followed only
by nonpositive differences through index $K(a)$.  The tail is nonincreasing
by Lemma~\ref{f60:lem-tail-cutoff}.  Hence the whole coefficient sequence
is unimodal.
\end{proof}

\subsection{The complete integer-row specification}
\label{f60:subsec-row-specification}

For clarity and reproducibility, the following specifies every row
family used by the finite linear certificates.  A row labelled $\rho$
means $\sum_{j,m}a^\rho_{j,m}t_{j,m}\le b^\rho$; unspecified
coefficients vanish.  All indices satisfy $1\le j\le v$ and
$0\le m\le a-j$.  The symbols in the rows are defined above, and
$r,h,l$ are integers.

\begin{description}
\item[\texttt{count}$(r)$.]
For $2\le r\le v$, use
$a_{j,m}=\mathbf{1}_{j=r}$ and $b=\binom vr$.

\item[\texttt{path}$(r)$.]
For $2\le r\le a$, use
$a_{j,m}=-\binom m{r-j}$ and
$b=\binom ar-\binom{n-r+1}{r}$.

\item[\texttt{private}$(r,h)$.]
For $2\le r\le v$ and $0\le h<a$, set
\[
 a_{j,m}=\mathbf{1}_{j=r}\bigl((r-1)m+a-r+1-rh\bigr)_+
 -(v-r+1)\mathbf{1}_{j=r-1}(m-h)_+,
 \qquad b=0.
\]

\item[\texttt{hall}$(r,l)$.]
Take $0\le l\le a$, $0\le r<\min(v,a-l)$, and
$c=v-r-\max(0,l-\delta)\ge0$.  Set
\[
 a_{j,m}=\bigl((r+1)\mathbf{1}_{j=r+1}-c\mathbf{1}_{j=r}\bigr)
                 \binom ml,\qquad
 b=\begin{cases}c\binom al,&r=0,\\0,&r>0.\end{cases}
\]
The implemented domain with $l=0$ uses $1\le r<v$; the omitted
$(r,l)=(0,0)$ case is the already imposed first-layer equality.

\item[\texttt{tail}$(r,h)$.]
For $2\le r\le v$ and $0\le h\le a-r+1$, set
\[
 a_{j,m}=\mathbf{1}_{j=r}b_{r,h}(m)
 -(v-r+1)\mathbf{1}_{j=r-1}\mathbf{1}_{m\ge h},\qquad b=0.
\]

\item[\texttt{release\_upper}$(r,h)$.]
On the same range, set
\[
 a_{j,m}=s_r\mathbf{1}_{j=r-1}(m-h)_+
            -\mathbf{1}_{j=r}U_{r,h}(m),\qquad b=0.
\]

\item[\texttt{tail\_upper}$(r,h)$.]
On the same range, set
\[
 a_{j,m}=s_r\mathbf{1}_{j=r-1}\mathbf{1}_{m\ge h}
            -\mathbf{1}_{j=r}V_{r,h}(m),\qquad b=0.
\]

\item[\texttt{mean}$(r)$.]
For $2\le r\le v$, set
\[
 a_{j,m}=-m\mathbf{1}_{j=r}-\beta_r\mathbf{1}_{j=2},\qquad
 b=-C_r^0-\beta_r\binom v2.
\]

\item[\texttt{union}$(r)$.]
For $2\le r\le v$, set
\[
 a_{j,m}=(a-m)\left(\mathbf{1}_{j=r}
                -\binom{v-1}{r-1}\mathbf{1}_{j=1}\right),\qquad b=0.
\]

\item[\texttt{edge\_lower}$(r)$.]
For $2\le r\le v$, set
\[
 a_{j,m}=-\mathbf{1}_{j=r}+\binom{v-2}{r-2}\mathbf{1}_{j=2},\qquad
 b=-\binom vr+\binom{v-2}{r-2}\binom v2.
\]

\item[\texttt{edge\_upper}$(r)$.]
For $2\le r\le v$ and $D_0>0$, set
\[
 a_{j,m}=D_0\mathbf{1}_{j=r}-G_r\mathbf{1}_{j=2},\qquad
 b=D_0\binom vr-G_r\binom v2.
\]
For $D_0=0$, use the \texttt{count}$(r)$ row instead.

\item[\texttt{assumption}$(k)$.]
Only in a \texttt{conditional} certificate at the same index $k$, use
\[
 a_{j,m}=B_{m,j}(k),\qquad b=-B_{a,0}(k).
\]
\end{description}

These formulas specify a relaxation containing the counts of every
forest after the prescribed choice of $I$.  They do not assert that
every point in the relaxation is realizable by a forest.

\subsection{An explicit affine comparison at the final parameter gap}
\label{f60:subsec-example}

For $(n,a,k)=(60,33,19)$ the data include a certificate with $193$
nonzero row multipliers, including one assumption row.  Its denominator
is $D_c=10^7$ and its equality-weight numerator is $1395$.
Retaining the assumption row algebraically on the left, rather than
declaring it true, gives the unconditional affine comparison
\begin{equation}
 \Delta_{20}\le\frac{96885639}{100000000}\Delta_{19}
 -\frac{4010142596414187852788421881}{28302876966000000000}.
 \label{f60:eq-example-exact}
\end{equation}
Since the last positive fraction is strictly greater than $141686748$,
\begin{equation}
 \Delta_{20}\le\frac{96885639}{100000000}\Delta_{19}-141686748.
 \label{f60:eq-example-rounded}
\end{equation}
The complete weights are in \texttt{example\_60\_33\_19.json}, which
agrees exactly with the corresponding entry of the master data file.
The comparison has been checked by recomputing both normalization
scales and the rational residual bound.  It applies simultaneously to
all forests with $(n,a)=(60,33)$, rather than to a selected graph.

\subsection{Exact obstructions to extending the stated relaxations}
\label{f60:subsec-countermodels}

The supplements also record two rational countermodels to specified
real-valued relaxations.  They are useful for delineating the reach of
the method, but have a different logical status from the finite forest
theorem.

The first, at $(n,a,k)=(55,30,18)$, satisfies all rows in the original
four families \texttt{count}, \texttt{path}, \texttt{private}, and
\texttt{hall} with $l\ge1$.  It has $193$ nonzero count variables and
a common denominator with $135$ decimal digits.  The second, at
$(100,55,31)$, satisfies the extended row families above, together
with $44$ additional valid cover rows
\begin{equation}
 T_r\ge\max\left(0,\binom vr-E\binom{v-2}{r-2}\right)
 \qquad(2\le r\le v).
 \label{f60:eq-cover}
\end{equation}
These follow from \eqref{f60:eq-edge-lower-raw} and $e\le E$.
The second countermodel has $693$ nonzero count variables and a common
denominator with $612$ decimal digits.  Its reverse-deletion rows with
$s_r=0$ need not be stored because they hold automatically for every
nonnegative array.

The checker \texttt{verify\_countermodels.py} regenerates the complete
specified row domains, rather than trusting a list of selected rows.
It checks nonnegativity, the first-layer equality, all inequalities,
and the coefficient differences using integer numerators and a common
positive denominator.  The resulting $1\,198$ and $8\,356$ row checks,
respectively, give
\begin{equation}
\begin{array}{c|cc}
 (n,a,k)&\Delta_k&\Delta_{k+1}\\
 \hline
 (55,30,18)&<-69\,000\,000&>74\,000\\
 (100,55,31)&<-2\cdot10^{18}&>3\cdot10^{17}.
\end{array}
 \label{f60:eq-countermodel-differences}
\end{equation}
Substitution into \eqref{f60:eq-decomposition} produces a formal
polynomial with nonnegative rational coefficients, constant coefficient
$1$, linear coefficient $n$, and degree $a$.  No graph realization of
either array is asserted.

It follows that the explicitly specified real-valued systems cannot
by themselves prove the no-recovery implication at these parameter
triples: a nonnegative linear combination of valid rows must remain
valid at a point satisfying all the rows.  This does not disprove
forest unimodality, rule out stronger consistency conditions or
integrality information, or rule out a different proof at order $100$.
It shows precisely why replacing the upper limit $60$ by $100$ in the
finite theorem is not justified by these certificates.


\section{The monotone window and the hard-core reduction}\label{sec:window}
This section fixes the interface used by the large-order chain. Its structural statements are independent of the finite-order relaxation.

\begin{lemma}[Initial and final monotone segments]\label{l59:window}
Let $F$ be an $n$-vertex forest with independence number $a$. Put
\[
 L=\lceil n/4\rceil,\qquad U=\lfloor(n+2a)/6\rfloor+1.
\]
Then $c_0\le\cdots\le c_L$ and $c_U\ge c_{U+1}\ge\cdots\ge c_a$.
\end{lemma}
\begin{proof}
Root every component. For an independent $k$-set $S$, consider triples $(S,v,u)$ with $v\in S$ and $u$ a child of $v$. If $u$ has no child in $S$, send the triple to $((S\setminus\{v\})\cup\{u\},u)$. Otherwise choose the first child $w$ of $u$ belonging to $S$ and send it to $(S,w)$. Outputs are independent $k$-sets with a marked nonroot vertex. In the first case the grandparent of the mark, if it exists, is absent; in the second it is the selected vertex $v$. The mark and its ancestors therefore recover the input, proving injectivity. If $D(S)=\sum_{v\in S}\deg(v)$ and $\rho_k=\sum_S|S\cap\{\text{roots}\}|$, this gives
\[
 \sum_S D(S)-kc_k+\rho_k\le kc_k-\rho_k.
\]
Counting extensions now yields $(k+1)c_{k+1}\ge(n-3k)c_k$. For every $k<L$, $n-3k\ge k+1$, proving the initial segment.

Fix a maximum matching of size $v=n-a$. Retaining only its edges gives a disjoint union of $v$ pairs and $a-v$ singletons, with polynomial
\[
 Q(x)=(1+2x)^v(1+x)^{a-v}=\sum q_kx^k.
\]
Uniform independent sets in two consecutive ranks of this reference graph admit an inclusion coupling. To see this, more generally give the blocks weights $w_1,\ldots,w_r>0$. Select a $k$-subset of blocks proportionally to the product of its weights. If $e_j$ denotes the elementary symmetric sum of the first $r-1$ weights, the probability of selecting the last block is $w_re_{k-1}/(e_k+w_re_{k-1})$. It increases with $k$, because the coefficient sequence of $\prod_{i<r}(1+w_ix)$ is log-concave. Couple these last-block indicators monotonically. If they agree, couple the remaining ranks by induction; if they differ, use the same remaining block set. Within common blocks use the same uniform element. All conditional marginals are preserved, including boundary ranks.

Independence in $F$ is a downward-closed event under this coupling, so $c_{k+1}/q_{k+1}\le c_k/q_k$. Finally
\[
 (k+1)(q_{k+1}-q_k)
 =(n+2a-6k+1)q_k-2(a-k+1)(q_k-q_{k-1}).
\]
For $k\ge U$ the first coefficient is nonpositive. If $q_k\ge q_{k-1}$ this proves $q_{k+1}\le q_k$ directly; otherwise log-concavity proves it. Ratio domination transfers this to $c_k$.
\end{proof}

At activity $\lambda>0$, sample an independent set with probability $\lambda^{|S|}/I_F(\lambda)$. Let $K=|S|$, $\mu=\E K$, and $V=\Var K$. We use the following scale-corrected mean lemma from Supplement A:
\begin{lemma}[Activity-two mean input]\label{l59:mean}
Every forest satisfies $3\mu_F(2)\ge a+29n/60$.
\end{lemma}
The complete proof of Lemma~\ref{l59:mean}, including its short-leg and terminal-branch reductions, is preserved among the foundational documents of stage 900. Its last step reduces a minimal negative tree to a leaf-rooted host with branches in
\[
 \{T_{33},T_{35},P_2,P_4,P_6,P_8,P_{10},P_{12}\}.
\]
The terminal rooted trees $T_{33},T_{35}$ have a root and two pendant paths of lengths $(3,3)$ and $(3,5)$. Each branch $T$ has root state $\alpha(T)-\alpha(T-\mathrm{root})=0$; this is an independence-number difference, not an occupation probability. With $G(F)=3\mu_F(2)-a-29n/60$, host ratio $r\in[5/3,3]$, $X=rG(H)$ and $Y=G(H-u)$, the source derives
\[
 E=60(Ar+B)G(F),\qquad E_2=180X+120Y+6r-54\ge0,
\]
where $A$ is the product of the branch partition functions with roots deleted and $A+B$ is the product without deletion. For $2\le m\le5$ branches, all 823 multisets containing a terminal tree satisfy
\[
 E-\frac A3E_2=A(60R-100)Y+A J(r)>0,\qquad R=(A+B)/A.
\]
Here $R\ge(7/5)^2$ and the least endpoint value of $J$ is $3523/473$. For all $m\ge6$, the remaining term is bounded below by
\[
 \left(\frac75\right)^m\left(\frac{343m}{1140}-\frac{89}{60}\right)
 -\frac{11m}{30}+\frac{89}{60}>0;
\]
its value at 6 is $5068593/2968750$ and successive increments are at least $11/60$. The earlier reductions prove that these cases exhaust every minimal negative tree; the 823-case arithmetic alone does not prove that exhaustion. The dependency map identifies all three source proofs and five reconstruction scripts.

Since $d\mu/d\log\lambda=V>0$ for a nonempty graph, the mean is strictly increasing. Conditioning on either bipartition class gives $V\ge\mu/[2(1+\lambda)]$, whence
\[
 \mu(9/4)\ge\sqrt{27/26}\,\mu(2)>(n+2a)/6.
\]
The last inequality follows from $a\ge n/2$ and $27\cdot59^2>26\cdot60^2$. Since every vertex has marginal occupation at most $\lambda/(1+\lambda)$, also $\mu(1/3)\le n/4$. Thus every necessary rank $L\le k<U$ has a unique activity $\lambda_k\in[1/3,9/4)$ with $\mu(\lambda_k)=k$.

For such an activity, let $\pi_v$ be the actual marginal at $v$, and choose an independent set $B$ maximizing $W=\sum_{v\in B}\pi_v$. This is a maximum-weight set, not necessarily a maximum-cardinality set. Bipartiteness gives $W\ge\mu/2$. Write $C=V(F)\setminus B$. Conditional on the entire occupied configuration in $C$, let $Y$ be its size and $M$ the number of unblocked vertices in $B$. Then
\begin{equation}\label{l59:mixture}
 K=Y+\operatorname{Bin}(M,q),\quad q=\frac\lambda{1+\lambda},\quad m=\E M,\quad qm=W.
\end{equation}
At $\mu=k$, put $j=k-Y$ and $\delta=j-qM$. The same probability space gives
\begin{equation}\label{l59:moments}
 \E\delta=0,\qquad \E\delta^2=V-q(1-q)m.
\end{equation}
No independence of $M$ and $Y$ is asserted.

\section{The final matching-count extension: orders 59 through 79}\label{sec:matching59}
We describe the part added by the outermost release. The older release proves the range $n\ge80$ using the same reduction.

\begin{lemma}[Leaves and the complement matching]\label{l59:matching}
Let $B$ be as in \eqref{l59:mixture}, $h=|B|$, and $H=F[C]$. Then
\[
 h\ge\lceil n/3\rceil,\quad \nu(H)\ge\max(0,n-2h),\quad \alpha(H)\le\lfloor n/2\rfloor.
\]
\end{lemma}
\begin{proof}
For a leaf $u$ whose component is not a single edge, let $v$ be its neighbor. The formulas
\[
 \pi_u=\frac{\lambda I_{F-\{u,v\}}(\lambda)}{I_F(\lambda)},\qquad
 \pi_v=\frac{\lambda I_{F-N[v]}(\lambda)}{I_F(\lambda)}
\]
show $\pi_u>\pi_v$, since the second induced graph is strictly smaller than the first. Replacing $v$ by $u$ would improve a maximum-weight independent set; therefore $v\notin B$, and maximality implies $u\in B$. Isolated vertices are in $B$, and in each edge component exactly one endpoint is in $B$.

Let $\kappa$ be the number of components of $F$, $d$ the number of its edge components, $c=n-h$, $e=|E(H)|$, and $s_0$ the number of isolated vertices in $H$. Each $C$-vertex has a neighbor in $B$. An isolated vertex of $H$ has at least two such neighbors, except for the $d$ endpoints in edge components. Hence $e(B,C)\ge c+s_0-d$, and $e+s_0\le h-\kappa+d$. The number of nontrivial components of $H$ is $c-e-s_0\ge n-2h$. Choosing one edge from each gives the matching bound. Comparing twice its size with $c$ gives $h\ge n/3$. If $h\le n/2$, then $\alpha(H)\le c-(n-2h)=h$; otherwise $\alpha(H)\le c=n-h$.
\end{proof}

For $c\ge2r$ define
\[
 D(c,r;y)=\sum_t\binom rt2^t\binom{c-2r}{y-t}.
\]
Deleting all edges of $H$ except a matching bounds its independent-set coefficients by those of $(1+2x)^r(1+x)^{c-2r}$. Therefore safe integer envelopes are
\begin{equation}\label{l59:envelope}
 A_{n,y}^{(M)}=\max_{\substack{\lceil n/3\rceil\le h\le n\\h\ge M}}
 D(n-h,\max(0,n-2h);y),\qquad A_{n,y}=A_{n,y}^{(0)}.
\end{equation}
An empty maximum is zero; binomial coefficients outside their natural ranges are zero.

\begin{lemma}[Exact cancellation]\label{l59:cancel}
For every integer $y$,
\[
 \E[\ind_{\{Y=y\}}(1-q)^M]
 =\frac{i_y(H)\lambda^y}{I_F(\lambda)}
 \le\frac{A_{n,y}\lambda^y}{P_n(\lambda)},
\]
where $P_n(x)=\sum_t\binom{n-t+1}{t}x^t$ is the path polynomial.
\end{lemma}
\begin{proof}
The probability of an independent configuration $S\subseteq C$ is $\lambda^{|S|}(1+\lambda)^{M(S)}/I_F(\lambda)$. Multiplication by $(1-q)^{M(S)}$ cancels the free-vertex factor. Sum over $|S|=y$ and use \eqref{l59:envelope} and the path coefficient lower bound proved in the finite-order part.
\end{proof}

Fix a rectangle $q\in[q_a,q_b]$, $m\in[m_l,m_h]$, and $59\le n\le79$. Put $a_0=q_a/(1-q_a)$, $b_0=q_b/(1-q_b)$. The certified density bound and integer rounding give $k\ge k_{\min}$ and $m\ge k_{\min}/(2q_b)$. All real target distributions are included in
\[
 \mathcal A=\{(n,k)\in\mathbb Z^2:59\le n\le79,\ k_{\min}\le k\le\lfloor2q_bm_h\rfloor,
 \ n/4\le k\le n/2,\ rn\le k\le q_bn\}.
\]
The value $r$ here is a certified density constant, not the matching size in $D$. Define fixed upward rational bounds
\[
 B_j\ge\max_{(n,k)\in\mathcal A}\frac{A_{n,k-j}b_0^{k-j}}{P_n(a_0)}.
\]
Terms with $k-j$ outside the range of independent-set sizes are zero. Lemma~\ref{l59:cancel} proves $\E[\ind_{\{k-Y=j\}}(1-q)^M]\le B_j$ throughout the rectangle. All possible conditional states belong to the finite set
\[
 \mathcal S=\{(M,j)\in\mathbb Z^2:0\le M\le79,\ \exists(n,k)\in\mathcal A:
 A_{n,k-j}^{(M)}>0\}.
\]
Excluding a state uses a zero combinatorial count, never a small numerical probability.

\section{Rational mixture certificates and the inherited chain}\label{sec:cert59}
Let $b_M(j;q)$ be the binomial point probability, extended by zero outside its support. Define
\[
 f_L=(1-q)b_M(j;q)-qb_M(j-1;q),\quad
 f_R=qb_M(j;q)-(1-q)b_M(j+1;q),
\]
\[
 f_C=2b_M(j;q)-b_M(j-1;q)-b_M(j+1;q).
\]
Every certificate uses fixed nonnegative weights $w_L,w_R,w_C$, not all zero, and $f=w_Lf_L+w_Rf_R+w_Cf_C$. Its scale $S$ is positive. Source inequalities on the same distribution provide
\begin{equation}\label{l59:sourcebounds}
 \E\delta^2\le\alpha_0m+\beta_0,\quad
 \Var M\le D_0m,\quad
 \E(1-q+qt_i)^M\le e^{-\ell_i m}.
\end{equation}
Here the finitely many $t_i\in[0,1]$ and $\ell_i\ge0$ are fixed rational parameters of the source bounds. Their proof is part of the chain, as described below; constants cannot be copied to a smaller threshold without checking the signs of the $n$-dependent terms.

The pointwise acceptance inequality is
\begin{equation}\label{l59:pointwise}
 Sf(M,j;q)\ge A_0+B_0M+C_0\delta-D_\delta\delta^2-D_MM^2
 -\sum_iT_i(1-q+qt_i)^M-U_j(1-q)^M,
\end{equation}
for every state in $\mathcal S$ and every $q$ in the rectangle, with $D_\delta,D_M,T_i,U_j\ge0$. The other coefficients may have either sign. Taking expectations gives
\begin{align}\label{l59:budget}
 S\E f\ge H(m):={}&A_0+B_0m-D_\delta(\alpha_0m+\beta_0)
 -D_M(m^2+D_0m)\notag\\
 &-\sum_iT_i e^{-\ell_i m}-\sum_jU_jB_j.
\end{align}
Since $H''(m)=-2D_M-\sum_iT_i\ell_i^2e^{-\ell_i m}\le0$, strictly positive endpoint budgets certify all real $m\in[m_l,m_h]$.

\subsection{Continuous activity certification}
For a fixed state let $R(q)$ be the left side of \eqref{l59:pointwise} minus the right side. Write $B_L(M)$, $B_R(M)$, and $B_C(M)$ for the supplied uniform upper bounds on the second $q$-derivatives of the three kernels, valid for every integer $j$ and all $q\in[q_a,q_b]$. They are reconstructed by \nolinkurl{derivative_bounds} in stage 80's \nolinkurl{bridge_compressed.py}, with their use specified in stage 59's proof, Section 9, and independent acceptance in \nolinkurl{check_matching_kernel.py}. These bounds give
\begin{align*}
 R''(q)\le Z_{M,j}:={}&S(w_LB_L(M)+w_RB_R(M)+w_CB_C(M))+2D_\delta M^2\\
 &+\sum_iT_iM(M-1)(1-t_i)^2(1-q_a+q_at_i)^{M-2}\\
 &+U_jM(M-1)(1-q_a)^{M-2}.
\end{align*}
The terms containing $M(M-1)$ are zero when $M<2$. With $h_q=q_b-q_a$, it suffices to check $R(q_a),R(q_b)\ge h_q^2Z_{M,j}/8$. The interpolation remainder then proves $R(q)\ge0$ on the entire interval. Rational enclosures for elementary functions and probabilities are directed; a precision scale is not a floating-point tolerance.

\subsection{The structural source interface}
For clarity, the frozen analytic sources use rooted-tree identities. If $p_v$ is the downward occupation message, $z_v$ its logarithmic response, and $r_v=\pi_v/p_v$, then
\[
 \frac{p_v}{1-p_v}=\lambda\prod_{u\text{ child of }v}(1-p_u),\quad
 z_v=1-\sum_up_uz_u,
\]
\[
 \mu=\sum_vp_vz_v=\sum_vp_vr_v,\qquad
 V=\sum_vp_vr_v(1-p_v)z_v^2.
\]
For $\lambda\le b$, $[1+b(1-p_v)]^{-1}\le r_v\le1$. Put
\[
 T_b(p)=\log\frac{b(1-p)}p,\quad \Lambda_b(p)=\frac1{pT_b(p)},\quad h(p)=1-p/2.
\]
If $0\le R_i\le b$ and $\sum_i\log(1+R_i)\le t$, convexity gives
\[
 \sum_iR_i\le H_b(t):=sb+e^{t-s\log(1+b)}-1,
 \qquad s=\lfloor t/\log(1+b)\rfloor.
\]
For the displayed scalar formulas take $0<p<q_b=b/(1+b)$. At the endpoint $p=q_b$, the vertex has no children and $z=1$; terms containing $(1-z)^2/T_b(p)$, $(1-z)_+^2/T_b(p)$, or $(z-1)_+^2/T_b(p)$ are defined to be zero there, as are the corresponding empty-child capacity terms. The checkers handle this leaf endpoint separately. Weighted Cauchy--Schwarz over disjoint child sets gives the summed nonpositive expressions
\[
 \sum_vp_v[\Lambda_b(p_v)(1-z_v)_+^2-h(p_v)(z_v)_+^2]\le0,
\]
\[
 \sum_vp_v[\Lambda_b(p_v)(z_v-1)_+^2-h(p_v)(z_v)_-^2]\le0,
\]
\[
 \sum_vp_v\left[\frac{r_v(1-z_v)^2}{H_b(T_b(p_v))}
 -(1-r_v)(1-p_v)z_v^2\right]\le0.
\]
These, additional message-flow identities, and $\sum_vp_v(z_v-r_v)=0$ transform the variance bounds into scalar inequalities in $(p,r,z)$. The response variable $z$ ranges over all real numbers. The checkers split its three intervals $(-\infty,0]$, $[0,1]$, and $[1,\infty)$ and prove polynomial nonnegativity, rather than sampling responses. Later stages replace $h(p)$ with the sharper $p/[-\log(1-p)]$ and add flow and logarithmic terms; the full formulas remain in their source documents.

For the Laplace bound in \eqref{l59:sourcebounds}, hold $B$ fixed and lower only its activities. The identity
\[
 \E e^{-s_\lambda(u)M}=\E e^{-uK_B},\qquad
 s_\lambda(u)=\log\frac{1+\lambda}{1+\lambda e^{-u}}
\]
allows integration of heterogeneous nonnegative-source variance bounds. Applying a common-activity variance bound to an unverified heterogeneous path would not suffice. The source proof explicitly establishes the needed heterogeneous domain.

The basic density source at activity 1 uses the path-mean result of Andriantiana, Razanajatovo Misanantenaina, and Wagner \cite{average}; it is not applied as a path extremality theorem at arbitrary activity. The activity-density improvements at stage 80 use 106 rational activity points. Their finite branch classes distinguish one-, two-, and three-vertex branches from larger ones; degrees 3 and 4 give 105 classes per point, while all higher degrees follow by removing one branch and an analytic induction. Thus the 11,130 checks are checks of the full induction scheme, not enumeration of forests up to some chosen order. Stage 59 adds four flow--log source certificates and recalculates the signed vertex costs with $59\le n\le79$.

\subsection{Coverage and exclusion of valleys}
The 231 new certificates cover 154 activity strata, subordinate to 43 parent activity intervals. Together with 188, 243, and 558 older finite kernels, respectively from stages 99, 130, and 170, they form 1220 finite rectangles. Their upper $m$ boundary joins the stage-350 unbounded-$m$ estimates. All 236 stage-80 kernels are excluded for $n<80$, because some of them use signed costs requiring $n\ge80$. This distinction is checked by the source and coverage code.

For the new certificates, the stored and replayed integer endpoint-state count is 1,185,408, and the smallest normalized expectation budget is
\[
 \frac{424830401847072141613691}{10^{28}}>0.
\]
The corresponding scale is 23; thus this normalized number must be divided by 23 to interpret it as a lower bound on the unscaled mixed-kernel expectation. It is not a common lower bound on a coefficient curvature.

Suppose at a necessary rank $k$ that $c_k\le c_{k-1}$ and $c_k\le c_{k+1}$. Writing $p_j=c_j\lambda^j/I_F(\lambda)$, the first two inequalities give $\E f_L\le0$ and $\E f_R\le0$. They also imply
\[
 p_{k-1}+p_{k+1}\ge(\lambda^{-1}+\lambda)p_k\ge2p_k,
\]
so $\E f_C\le0$. This contradicts the strictly positive budget in \eqref{l59:budget}. Every necessary rank therefore satisfies $c_k>\min(c_{k-1},c_{k+1})$. If a strict descent is followed later by a strict ascent, choose the first subsequent index $k$ with $c_{k+1}\ge c_k$. Then $c_{k-1}>c_k$ and $c_{k+1}\ge c_k$. The monotone end segments in Lemma~\ref{l59:window} place such an index in $[L,U-1]$, giving a prohibited weak local minimum, including the possibility of a flat valley. Consequently all orders $59\le n\le79$ are covered once the source and coverage certificates are accepted.

\begin{table}[ht]
\centering
\caption{Frozen large-order stages. Each row proves the new finite interval and inherits the larger orders; the base stage includes an unbounded-order argument.}
\begin{tabular}{rrl}
\toprule Stage & New orders & Mathematical source in Supplement A\\\midrule
900 & $n\ge900$ & base source and kernel proof\\
500 & $500\le n\le899$ & stage 500 proof and replay\\
350 & $350\le n\le499$ & stage 350 proof and replay\\
170 & $170\le n\le349$ & stage 170 proof and replay\\
130 & $130\le n\le169$ & stage 130 proof and replay\\
99 & $99\le n\le129$ & stage 99 proof and replay\\
80 & $80\le n\le98$ & stage 80 proof and replay\\
59 & $59\le n\le79$ & final matching-count extension\\\bottomrule
\end{tabular}
\end{table}
The recursive replay authenticates and checks these frozen stages. The original mathematical proofs are included to justify their all-order and continuous-domain claims. Taking the base and each extension in the table establishes Proposition~\ref{prop:large}; this is a dependency-chain proof, not an extrapolation from the 79-vertex conditional-state bound.


\section{Reproducibility, dependencies, and interpretation}\label{sec:repro}
The mathematical trust boundary has three parts. First, structural proofs must show that every forest satisfies the asserted inequalities. Second, arithmetic checkers must implement those inequalities with safe rounding and verify the full parameter domains. Third, the domains must cover all required orders, ranks, activities, and conditional states. Passing one of these checks is not a substitute for the others.

The finite-order verifier constructs the full scalar domain instead of trusting a list of covered parameters. It rejects missing, duplicate, and extraneous stored certificates, inadmissible rows, and negative inequality multipliers. The large-order checkers independently reconstruct matching-count envelopes and the discrete $(n,k)$ and $(M,j)$ domains; continuous activity intervals and real $m=\E M$ intervals are certified by derivative or concavity arguments.

\subsection{Commands and preserved inputs}
The wrapper \texttt{verification/reproduce.py} in the distribution checks the input archive hashes, extracts into a user-selected working directory, runs the finite-order and countermodel verifiers, and, in full mode, executes the complete large-order chain. The underlying commands are:
\begin{verbatim}
python -S verify.py
python -S verify_grid.py
python -S verify_countermodels.py
python -S check_release.py
python -S replay.py
\end{verbatim}
The first three run in Supplement B; the last two run in Supplement A. Use Python 3.10 or later and do not use \texttt{-O}. The computational proof acceptance path needs only the standard library. Optional discovery routines and numerical optimizers are not called by acceptance. The large-order chain can be much slower than the finite-order part and creates nested working directories; a short extraction path is helpful on Windows.

Fresh runs, inherited records, and diagnostic checks are separated in the accompanying audit report. A run with \texttt{--skip-prior} authenticates inherited bytes and checks the new segment, but is not described as rerunning all earlier stages. The source archive hashes and final distribution manifest identify the exact mathematical data. Runtime-dependent audit files are not frozen mathematical premises.

\subsection{Precise limits of the 100-vertex obstruction}
The array in \texttt{countermodel\_100\_55\_31.json} satisfies the stated nonnegative-real relaxation for $(n,a,k)=(100,55,31)$, including its extra covering rows, and has $\Delta_{31}<0<\Delta_{32}$. This proves that those relaxation constraints alone cannot imply the needed no-valley conclusion. The formal polynomial is not asserted to be the independence polynomial of any graph. Proposition~\ref{prop:large} uses additional probabilistic and structural information, so there is no contradiction between that obstruction and the consolidated assertion. The analogous order-55 array concerns the older, weaker finite-order model.

\subsection{Computational and editorial assistance}
AI-assisted tools were used for mathematical review, exact computational reproduction, literature research, translation, and manuscript preparation. The resulting audit reports document their scope and limitations; they are not independent human referee reports or proof-assistant certificates. All mathematical claims are submitted for review together with the written arguments, exact certificates, and complete dependency sources. Submission instructions and provenance records are provided separately from the mathematical text.

\begin{thebibliography}{9}
\bibitem{alavi} Y. Alavi, P. J. Malde, A. J. Schwenk, and P. Erd\H{o}s, The vertex independence sequence of a graph is not constrained, \emph{Congressus Numerantium} \textbf{58} (1987), 15--23.
\bibitem{erdos} Erd\H{o}s Problems, Problem 993, \url{https://www.erdosproblems.com/993}, accessed 26 September 2026.
\bibitem{jsp} The Justin Sun Prize, Problem JSP-000826 and contribution guidelines, \url{https://github.com/TheJustinSunPrize/awards/blob/main/problems/catalog-0801-0900.md}, \url{https://github.com/TheJustinSunPrize/awards/blob/main/CONTRIBUTING.md}, accessed 26 September 2026.
\bibitem{fang} E. X. Fang, J. Lu, E. Nevo, Y. Yao, and H. Zheng, Unimodality of Independence Polynomials for Sufficiently Large Forests, arXiv:2609.20961v1 (2026), \url{https://arxiv.org/abs/2609.20961v1}.
\bibitem{average} E. O. D. Andriantiana, V. Razanajatovo Misanantenaina, and S. Wagner, The average size of independent sets of graphs, arXiv:1807.08290, \url{https://arxiv.org/abs/1807.08290}.
\bibitem{kadrawi} O. Kadrawi and V. E. Levit, The independence polynomial of trees is not always log-concave starting from order 26, \emph{Ars Mathematica Contemporanea} \textbf{25} (2025), P4.03, \url{https://doi.org/10.26493/1855-3974.3207.2ad}.
\bibitem{ramos} E. Ramos and S. Sun, An AI enhanced approach to the tree unimodality conjecture, arXiv:2510.18826 (2025), \url{https://arxiv.org/abs/2510.18826}.
\end{thebibliography}
\end{document}
